\documentclass[11pt,a4paper]{article}
\usepackage{ctex}
\usepackage{graphicx}
\usepackage{geometry}
\usepackage{listings}
\usepackage{xcolor}
\usepackage{float}
\usepackage{enumitem}

\geometry{margin=2.5cm}

\lstset{
    language=bash,
    basicstyle=\ttfamily\small,
    backgroundcolor=\color{gray!5},
    frame=single,
    breaklines=true,
    numbers=none,
    showstringspaces=false
}

\title{系统开发工具实验报告2}
\author{姓名：李博文 \qquad 学号：25020007057}
\date{}

\begin{document}
\maketitle

\section*{第5题 控制一个可清理的后台任务}
\textbf{主题：}命令行环境：进程、信号与任务控制 \quad 建议用时：12—15分钟

\begin{enumerate}
\item 创建并进入 q05 目录
\begin{lstlisting}
mkdir -p ~/q05
cd ~/q05
\end{lstlisting}

\item 编辑 worker.sh，写入题目给定全部脚本内容
\begin{lstlisting}
nano worker.sh
\end{lstlisting}
worker.sh脚本内容：
\begin{lstlisting}
#!/usr/bin/env bash
trap 'echo CLEAN_EXIT >> cleanup.log; exit 0' TERM INT
n=0
while true; do
echo "$n"
n=$((n+1))
sleep 1
done
\end{lstlisting}

\item 给脚本添加执行权限
\begin{lstlisting}
chmod +x worker.sh
\end{lstlisting}

\item 前台启动脚本，标准输出、标准错误分别重定向日志文件
\begin{lstlisting}
./worker.sh > stdout.log 2> stderr.log
\end{lstlisting}
此时终端被脚本占用，不会在屏幕打印数字，数字全部写到 stdout.log。

\item 按 \texttt{Ctrl+Z} 挂起当前前台任务

\item 使用 \texttt{bg} 命令，把挂起的任务放到后台继续运行

\item 用 \texttt{jobs} 查看后台任务，确认状态为 Running
\begin{lstlisting}
jobs
\end{lstlisting}

\item 通过 \texttt{jobs -p} 获取任务 PID，发送 SIGTERM 信号，禁止 kill -9
\begin{lstlisting}
kill $(jobs -p)
\end{lstlisting}

\item 执行 jobs 确认任务已经退出，jobs 输出为空
\begin{lstlisting}
jobs
\end{lstlisting}

\item 检查清理日志，验证是否输出CLEAN\_EXIT
\begin{lstlisting}
cat cleanup.log
\end{lstlisting}
\end{enumerate}

\begin{figure}[H]
    \centering
    \includegraphics[width=0.9\textwidth]{img/q05.png}
    \caption{q05实验运行截图}
\end{figure}


\section*{第6题 语义重构与本地开发反馈}
\textbf{主题：}开发环境与工具 \quad 建议用时：12—15分钟
\begin{enumerate}
\item 先安装 pipx
\begin{lstlisting}
sudo apt update
sudo apt install -y pipx
\end{lstlisting}

\item 用 pipx 安装 ruff、pytest
\begin{lstlisting}
pipx install ruff
pipx install pytest
\end{lstlisting}

\item 把 pipx 的程序路径加入 shell
\begin{lstlisting}
pipx ensurepath
\end{lstlisting}

\item 创建三个 py 文件：
\begin{lstlisting}
cat > math_utils.py <<'EOF'
def total_price(price: float, count: int)-> float:
    return price * count
EOF

cat > app.py <<'EOF'
from math_utils import total_price
print(total_price(12.5, 4))
EOF

cat > test_math_utils.py <<'EOF'
from math_utils import total_price
def test_total_price():
    assert total_price(12.5, 4) == 50.0
EOF
\end{lstlisting}

\item 打开Vs code，连上 WSL Ubuntu
\begin{itemize}
\item 打开 app.py，光标放到 \texttt{total\_price}，按 F12，跳转到函数定义。
\begin{figure}[H]
    \centering
    \includegraphics[width=0.85\textwidth]{img/q061.png}
    \caption{F12跳转函数定义截图}
\end{figure}

\item 光标放在函数名\texttt{total\_price}，按 \texttt{Shift+F12}，查看所有引用位置。
\begin{figure}[H]
    \centering
    \includegraphics[width=0.85\textwidth]{img/q062.png}
    \caption{Shift+F12查看全部引用截图}
\end{figure}

\item 光标放在函数名，按 \texttt{F2}，输入新名字：\texttt{calculate\_total}，回车。
\begin{figure}[H]
    \centering
    \includegraphics[width=0.85\textwidth]{img/q063.png}
    \caption{F2语义重命名截图}
\end{figure}
\end{itemize}

\item 回到WSL终端，执行检查与测试
\begin{lstlisting}
ruff check .
ruff check . --fix
pytest
\end{lstlisting}
\end{enumerate}

\begin{figure}[H]
    \centering
    \includegraphics[width=0.9\textwidth]{img/q064.png}
    \caption{ruff与pytest运行截图}
\end{figure}


\section*{第7题 用调试器定位归并排序缺陷}
\textbf{主题：}Debugging \quad 建议用时：12—15分钟
\begin{enumerate}
\item 创建目录与文件
\begin{lstlisting}
mkdir -p ~/q07
cd ~/q07
\end{lstlisting}

\item 新建merge\_sort.py，复制下面带 bug 原始代码
\begin{lstlisting}
cat > merge_sort.py <<'EOF'
def merge(left, right):
    import pdb; pdb.set_trace()
    result = []
    i = j = 0
    while i < len(left) and j < len(right):
        if left[i] <= right[j]:
            result.append(left[i]); i += 1
        else:
            result.append(right[i]); j += 1
    return result + left[i:] + right[j:]

def merge_sort(a):
    if len(a) <= 1:
        return a
    m = len(a) // 2
    return merge(merge_sort(a[:m]), merge_sort(a[m:]))

print(merge_sort([3, 1, 4, 1, 5, 9, 2, 6]))
EOF
\end{lstlisting}

\item 运行原始有 bug 代码
\begin{lstlisting}
python3 merge_sort.py
\end{lstlisting}

\begin{figure}[H]
    \centering
    \includegraphics[width=0.9\textwidth]{img/q071.png}
    \caption{bug版本运行截图}
\end{figure}

\item pdb 分步调试

重复输入 n 多次，直到执行进入 \texttt{while i < len(left) and j < len(right):} 这一行。输入：\texttt{p i,j,left,right}。继续重复输入n多次，当程序跑到 else 那一行：\texttt{result.append(right[i]); j += 1}，输入\texttt{p i,j}，输入\texttt{q}退出调试。

\begin{figure}[H]
    \centering
    \includegraphics[width=0.9\textwidth]{img/q072.png}
    \caption{pdb调试过程截图}
\end{figure}

\item 写入修复完成的代码（删掉 pdb 断点，\texttt{right[i]}改成\texttt{right[j]}）
\begin{lstlisting}
cat > merge_sort.py <<'EOF'
def merge(left, right):
    result = []
    i = j = 0
    while i < len(left) and j < len(right):
        if left[i] <= right[j]:
            result.append(left[i]); i += 1
        else:
            result.append(right[j]); j += 1
    return result + left[i:] + right[j:]

def merge_sort(a):
    if len(a) <= 1:
        return a
    m = len(a) // 2
    return merge(merge_sort(a[:m]), merge_sort(a[m:]))

print(merge_sort([3, 1, 4, 1, 5, 9, 2, 6]))
EOF
\end{lstlisting}

\item 运行修复后的程序
\begin{lstlisting}
python3 merge_sort.py
\end{lstlisting}

\begin{figure}[H]
    \centering
    \includegraphics[width=0.9\textwidth]{img/q073.png}
    \caption{修复后运行截图}
\end{figure}

\item 创建 pytest 测试文件
\begin{lstlisting}
cat > test_merge.py <<'EOF'
from merge_sort import merge_sort

def test_given_input():
    assert merge_sort([3, 1, 4, 1, 5, 9, 2, 6]) == [1,1,2,3,4,5,6,9]

def test_dup_list():
    assert merge_sort([7,2,2,5,2]) == [2,2,2,5,7]
EOF
\end{lstlisting}

\item 执行测试
\begin{lstlisting}
pytest
\end{lstlisting}
\end{enumerate}

\begin{figure}[H]
    \centering
    \includegraphics[width=0.9\textwidth]{img/q074.png}
    \caption{pytest全部用例通过截图}
\end{figure}


\section*{第8题 先测量，再优化慢速词频程序}
\textbf{主题：}Profiling \quad 建议用时：12—15分钟
\begin{enumerate}
\item 创建目录，生成文件
\begin{lstlisting}
mkdir -p ~/q08
cd ~/q08
\end{lstlisting}

一． generate\_words.py 生成词表
\begin{lstlisting}
cat > generate_words.py <<'EOF'
import random
random.seed(20260831)
vocab = [f"w{i}" for i in range(1000)]
with open("words.txt", "w", encoding="utf-8") as f:
    f.write(" ".join(random.choice(vocab) for _ in range(30000)))
EOF
\end{lstlisting}
运行生成words.txt
\begin{lstlisting}
python3 generate_words.py
\end{lstlisting}

二． 原始程序 wordfreq.py
\begin{lstlisting}
cat > wordfreq.py <<'EOF'
words = open("words.txt", encoding="utf-8").read().split()
unique = []
for word in words:
    if word not in unique:
        unique.append(word)
print("count=", len(unique))
EOF
\end{lstlisting}

\item time 测耗时，运行 2 次
\begin{lstlisting}
time python3 wordfreq.py
time python3 wordfreq.py
\end{lstlisting}
\end{enumerate}

\begin{figure}[H]
    \centering
    \includegraphics[width=0.9\textwidth]{img/q081.png}
    \caption{原始版本time计时截图}
\end{figure}

\begin{enumerate}[resume]
\item cProfile 找耗时点
\begin{lstlisting}
python3 -m cProfile -s cumulative wordfreq.py
\end{lstlisting}
\end{enumerate}

\begin{figure}[H]
    \centering
    \includegraphics[width=0.9\textwidth]{img/q082.png}
    \caption{cProfile性能分析截图}
\end{figure}

\begin{enumerate}[resume]
\item 优化版本 wordfreq\_opt.py
\begin{lstlisting}
cat > wordfreq_opt.py <<'EOF'
words = open("words.txt", encoding="utf-8").read().split()
unique_set = set()
for word in words:
    unique_set.add(word)
unique = list(unique_set)
print("count=", len(unique))
EOF
\end{lstlisting}

\item 运行 2 次计时
\begin{lstlisting}
time python3 wordfreq_opt.py
time python3 wordfreq_opt.py
\end{lstlisting}
\end{enumerate}

\begin{figure}[H]
    \centering
    \includegraphics[width=0.9\textwidth]{img/q083.png}
    \caption{优化后版本计时截图}
\end{figure}


\section*{练习}
\subsection*{练习1}
\textbf{题目：}你可能见过像 \texttt{cmd --flag -- --notaflag} 这样的命令。这里的 \texttt{--} 是一个特殊参数，它告诉程序后面不要再继续解析选项（flag）了。也就是说，\texttt{--} 后面的所有内容都会被当作位置参数（positional argument）。这有什么用？试着运行 \texttt{touch -- -myfile}，然后在不使用 \texttt{--} 的情况下把它删掉。

\textbf{答：}
\texttt{--}是命令行选项终止符，通知程序停止解析后续以\texttt{-}开头的内容为命令选项 (flag)，\texttt{--}之后全部内容视作普通位置参数。

\begin{figure}[H]
    \centering
    \includegraphics[width=0.85\textwidth]{img/s1.png}
    \caption{练习1运行截图}
\end{figure}

\subsection*{练习2}
\textbf{题目：}进程替换 \texttt{<(command)} 可以让你把一个命令的输出当成文件来用。试着配合 \texttt{diff} 和进程替换，比较 \texttt{printenv} 与 \texttt{export} 的输出。它们为什么不一样？（提示：可以试试 \texttt{diff <(printenv | sort) <(export | sort)}）

\textbf{答：}
\begin{lstlisting}
diff <(printenv | sort) <(export | sort)
\end{lstlisting}
1.输出格式不同。\texttt{export} 输出是 shell 声明语句：\texttt{declare -x HOME="/home/xxx"}；\texttt{printenv} 只输出变量=值的原始内容，没有 declare -x 前缀。即使变量集合相同，文本行不一样，diff 会报很多行差异。

2.变量集合不完全等同。
\texttt{printenv}：全部已经导出生效的环境变量，会包含继承来的环境变量。
\texttt{export} 只列出 shell 里被标记为要导出的变量；有些 shell 内部变量不一定出现在两者中。
被 export 标记的变量才会进入 printenv 的环境；但二者输出文本格式完全不同，排序后 diff 仍然看到大量不同。

\subsection*{练习3}
\textbf{题目：}阅读 \texttt{man ls}，写出一个 ls 命令，让它按下面的方式列出文件：
\begin{itemize}
\item 包含所有文件，包括隐藏文件
\item 文件大小以易读格式显示（例如 454M，而不是 454279954）
\item 按最近修改时间排序
\item 输出带颜色
\end{itemize}

示例输出大概像这样：
\begin{verbatim}
 -rw-r--r--   1 user group 1.1M Jan 14 09:53 baz
 drwxr-xr-x   5 user group  160 Jan 14 09:53 .
 -rw-r--r--   1 user group  514 Jan 14 06:42 bar
 -rw-r--r--   1 user group 106M Jan 13 12:12 foo
 drwx------+ 47 user group 1.5K Jan 12 18:08 ..
\end{verbatim}

\textbf{答：}
\begin{lstlisting}
ls -lath --color
\end{lstlisting}

\begin{figure}[H]
    \centering
    \includegraphics[width=0.9\textwidth]{img/s3.png}
    \caption{练习3 ls命令截图}
\end{figure}

\subsection*{练习4}
\textbf{题目：}写两个 bash 函数 \texttt{marco} 和 \texttt{polo}，行为如下：每次执行 \texttt{marco} 时，都要以某种方式保存当前工作目录；之后无论你切到哪个目录，只要执行 \texttt{polo}，它都应该把你 cd 回执行 marco 时所在的目录。为了方便调试，你可以把代码写进 \texttt{marco.sh}，然后执行 \texttt{source marco.sh} 把这些定义重新加载到当前 shell。

\textbf{答：}
marco.sh代码：
\begin{lstlisting}
marco() {
    MARCO_SAVED_DIR="$PWD"
}

polo() {
    cd "$MARCO_SAVED_DIR"
}
\end{lstlisting}
加载脚本到当前shell：
\begin{lstlisting}
source marco.sh
\end{lstlisting}

\begin{figure}[H]
    \centering
    \includegraphics[width=0.9\textwidth]{img/s4.png}
    \caption{练习4 marco polo截图}
\end{figure}

\subsection*{练习5}
\textbf{题目：}创建一个 alias dc，让你把 cd 打错时也能正常工作。

\textbf{答：}
\begin{lstlisting}
alias dc='cd'
\end{lstlisting}

\begin{figure}[H]
    \centering
    \includegraphics[width=0.6\textwidth]{img/s5.png}
    \caption{练习5 alias dc截图}
\end{figure}

\subsection*{练习6}
\textbf{题目：}运行 \texttt{history | awk '{\$1="";print substr(\$0,2)}' | sort | uniq -c | sort -n | tail -n 10}，找出你最常用的 10 条命令，并考虑给它们写更短的 alias。注意：这条命令适用于 Bash；如果你用的是 ZSH，把 \texttt{history} 换成 \texttt{history 1}。

\textbf{答：}
\begin{lstlisting}
history | awk '{\$1="";print substr(\$0,2)}' | sort | uniq -c | sort -n | tail -n 10
\end{lstlisting}
>提示zsh用户将history替换为\texttt{history 1}。

\begin{figure}[H]
    \centering
    \includegraphics[width=0.9\textwidth]{img/s6.png}
    \caption{练习6历史命令统计截图}
\end{figure}

\subsection*{练习7}
\textbf{题目：}调试排序算法：以下伪代码实现了合并排序，但存在一个错误。用你选择的语言实现它，然后用调试器（gdb、lldb、pdb，或你IDE的调试器）来查找并修复bug。
\begin{verbatim}
function merge_sort(arr):
    if length(arr) <= 1:
        return arr
    mid = length(arr) / 2
    left = merge_sort(arr[0..mid])
    right = merge_sort(arr[mid..end])
    return merge(left, right)

function merge(left, right):
    result = []
    i = 0, j = 0
    while i < length(left) AND j < length(right):
        if left[i] <= right[j]:
            append result, left[i]
            i = i + 1
        else:
            append result, right[i]
            j = j + 1
    append remaining elements from left and right
    return result
\end{verbatim}
测试向量：\texttt{merge\_sort([3, 1, 4, 1, 5, 9, 2, 6])}，应返回\texttt{[1, 1, 2, 3, 4, 5, 6, 9]}。使用断点并逐步通过合并函数找出错误元素被选中的位置。

\textbf{答：}
错误位置：merge函数内else分支写 \texttt{result.append(right[i]); j +=1}，下标错误，应该是\texttt{right[j]}。

\begin{figure}[H]
    \centering
    \includegraphics[width=0.85\textwidth]{img/s71.png}
    \caption{练习7pdb调试截图1}
\end{figure}
\begin{figure}[H]
    \centering
    \includegraphics[width=0.85\textwidth]{img/s72.png}
    \caption{练习7pdb调试截图2}
\end{figure}

\subsection*{练习8}
\textbf{题目：}使用 AddressSanitizer 调试内存错误。保存为：\texttt{uaf.c}
\begin{lstlisting}
#include <stdlib.h>
#include <string.h>
#include <stdio.h>

int main() {
    char *greeting = malloc(32);
    strcpy(greeting, "Hello, world!");
    printf("%s\n", greeting);

    free(greeting);

    greeting[0] = 'J';
    printf("%s\n", greeting);

    return 0;
}
\end{lstlisting}

首先编译并运行时不使用消毒剂：\texttt{gcc uaf.c -o uaf \&\& ./uaf}。看起来它有效。现在用 AddressSanitizer 编译：\texttt{gcc -fsanitize=address -g uaf.c -o uaf \&\& ./uaf}。阅读错误报告。ASan发现了什么bug？解决它识别出的问题。

\textbf{答：}
编译运行命令：
\begin{lstlisting}
gcc uaf.c -o uaf && ./uaf
gcc -fsanitize=address -g uaf.c -o uaf && ./uaf
\end{lstlisting}

ASan检测出 \textbf{heap‑use‑after‑free（释放后继续访问堆内存）}。内存free之后不允许读写这块内存。修复：调整free放到读写之后。

\begin{figure}[H]
    \centering
    \includegraphics[width=0.9\textwidth]{img/s81.png}
    \caption{练习8 ASan报错截图}
\end{figure}
\begin{figure}[H]
    \centering
    \includegraphics[width=0.9\textwidth]{img/s82.png}
    \caption{练习8修复代码运行截图}
\end{figure}

\subsection*{练习9}
\textbf{题目：}用来获取你选择的程序的基本性能统计数据。不同的计数器是什么意思？

\textbf{答：}
\begin{itemize}
\item \texttt{cycles}:CPU时钟周期；\texttt{instructions}执行指令，二者比值 IPC 衡量 CPU 流水线效率。
\item \texttt{branches / branch‑misses}统计分支指令与预测失败，分支失败会造成性能损耗。
\item \texttt{task‑clock}程序占用 CPU 时间；\texttt{context‑switches}上下文切换；\texttt{cpu‑migrations}进程跨核心迁移；\texttt{page‑faults}虚拟内存缺页。
\item \texttt{elapsed}是程序真实总运行时间。
\end{itemize}

\subsection*{练习10}
\textbf{题目：}用于在运行资源密集型程序时监控系统。尝试用 htop、taskset 来限制进程可以使用的 CPU：\texttt{taskset --cpu-list 0,2 stress -c 3}。为什么不使用三颗CPU？

\textbf{答：}
命令：
\begin{lstlisting}
taskset --cpu-list 0,2 stress -c 3
\end{lstlisting}
原因：\texttt{taskset --cpu-list}用来限制进程允许调度的 CPU 核心；\texttt{stress -c N}产生 N 个 CPU 压力线程。命令 \texttt{taskset --cpu-list 0,2 stress -c 3}：CPU 列表只提供核心 0、2，仅 2 颗核心，却启动 3 个 CPU 密集线程。 因此无法使用 3 颗 CPU；3 个线程竞争 2 个 CPU，会出现大量上下文切换，CPU1 始终空闲。 如果要使用 3 颗 CPU，cpulist 应当写 3 个不同核心编号，例如\texttt{0,1,2}。

\end{document}